
This study set out to measure which Pile domains drive which NLP benchmarks. Two foundational results are confirmed with high confidence: The Pile's domains are measurably distinct in cognitive content ($\eta^2$ > 0.91 for entity density, narrative coherence, and formal syntax density), and within-family domain exposure trajectories are genuinely non-uniform (10 of 22 domains with std > 0.001 across 154 Pythia checkpoints, confirmed to be scale-independent). These two results establish the empirical prerequisites for domain-benchmark specificity analysis.

The primary specificity test --- panel OLS regression mapping domain exposure to benchmark-specific improvement --- was blocked by a structural data gap: Books3 domain exposure is zero in the 600,000-document first-shard sample, and the benchmark evaluation cache was incomplete at analysis time. The panel regression pipeline (2,943 lines, 21/24 unit tests passing) is fully implemented and validated, but could not be executed. Preliminary Spearman analysis at 70M scale (N = 10 checkpoints) shows a direction reversal relative to prediction, most plausibly attributable to scale floor effects and sampling artifact rather than a definitive conclusion about domain-benchmark relationships.

The domain-benchmark specificity hypothesis is not resolved but is precisely scoped: the required data conditions (full 134M-document lookup, complete evaluation cache for $\geq 3$ model sizes) are characterized, and the execution pathway is defined. The infrastructure contribution --- validated trajectory extraction from Pythia's exact dataloaders, the panel regression pipeline, and documentation of the first-shard coverage limitation --- provides a reusable foundation for future domain attribution studies.

